\section{Methods}
\label{sec:methods}

\subsection{Model}
\label{sec:model}

The model description follows the ODD protocol \citep{Grimm2020}; the full ODD and the code are
in the Supplementary Material (SM, Section~\ref{SM-supp:odd}). The model is individual based and stochastic, with
one time step per week. Individuals live in a closed arena of $|G|=L\times L$ cells ($L=30$). A cell stands for a
nest site or local group: it is the unit at which crowding is experienced and mates are met. Each individual has an
age $a$, a sex, a position and a \emph{reproductive competence} $s\in[0,1]$, a phenomenological scalar that summarises
the capacity to breed, the capacity to rear young and the quality of the individual as a demonstrator for juveniles.
There are no resources and no predators, and no individual knows the population size. The simple interaction rules
of the title are all the rules below: the core processes, among them crowding damage and senescence, which suffice
for the collapse (Section~\ref{sec:res-wm}), and four social rules, which add specific patterns.
Table~\ref{tab:params} lists the key parameters.

\paragraph{Core processes}
\emph{Crowding damage.} Each week, after moving, competence is updated with the occupancy $m_i$ of the cell of $i$,
itself included:
\begin{equation}\label{eq:stress}
  s_i\leftarrow\min\{1,\max[0,\;s_i+g_i-\gamma\,(m_i-m_c)_+]\},
\end{equation}
where $g_i$ is the weekly gain (below), $m_c$ the tolerated group size and $\gamma$ the damage per extra occupant. In
Universe~25 social disruption arose while food and water were in excess \citep{Calhoun1973}, so the rule is
interference without a resource. It responds to the crowding that an individual experiences, not to the density
$N/|G|$.
\emph{Competence gates breeding and care.} A reproductive female that shares her cell with a reproductive male
(chosen at random) conceives with probability $p_0s_fs_m$; each pup survives birth with probability $q_0s_f$ and
starts with a fraction $w$ of its mother's competence,
\begin{equation}\label{eq:breed}
  p_{\rm conc}=p_0\,s_f\,s_m,\qquad p_{\rm surv}=q_0\,s_f,\qquad s_{\rm pup}=w\,s_f.
\end{equation}
The last relation is a maternal effect \citep{GinzburgTaneyhill1994}, not inheritance; it is an
assumption of \Cone, not an ingredient of the mechanism ($w=0$ with faster learning gives the same outcome,
Section~\ref{sec:res-patterns}). The rule encodes the failures that Calhoun reported in the crowded phase: males that
ceased courting, and mothers that abandoned their litters.
\emph{Senescence and attraction.} Death is senescent, with weekly hazard $[1+e^{-k(a-a_{\max})}]^{-1}$, independent
of $s$; females breed between $a_{\min}$ and $a_{\rm rep}$ weeks of age. These values follow Universe~25: menopause at
about $560$ days, many mice alive at about $800$ days, and withdrawn animals that were ``physically healthy''
\citep{Calhoun1973}. Each week every individual moves to a cell $y$ of its $3\times3$ neighbourhood (own cell
included) with probability proportional to
\begin{equation}\label{eq:move}
  W_i(y)=\bigl[1+\alpha\,\chi_i(y)\bigr]\;e^{-\kappa(1-s_i)^{p}n_y}\;\bigl[1+\pi(1-s_i)^{q}\,\mathbf 1_{\mathcal R}(y)\bigr],
\end{equation}
where $n_y$ is the number of other individuals in $y$ and $\chi_i(y)$ the summed affinity of $i$ for them; affinities
grow by proximity and decay with a memory of about 100 weeks. The first factor, attraction to acquaintances,
aggregates the colony, in the spirit of Calhoun's ``behavioural sink'' \citep{Calhoun1962}.

\paragraph{Four social rules}
The four social rules are labelled as in the code and in the figures.
\emph{Sensitive window (R1).} Competence can be gained only up to age $a_w$, a window that spans weaning and
maturity ($g_i=0$ for $a_i>a_w$), while crowding damage acts at every age. Mice removed from a declining companion
colony bred very poorly in uncrowded enclosures: after two months without reproduction, half of the females
conceived with normal males and 5 of about 100 pups survived \citep{Marsden1972}. The rule is the standard form of a
developmental sensitive window \citep{FawcettFrankenhuis2015}. Of all the assumptions, adult irreversibility has the
least empirical support and the largest effect; a variant multiplies the adult gain by a factor $\eadult\in[0,1]$,
the adult recovery factor (\Cone: $\eadult=0$).
\emph{Social learning of competence (R2).} Inside the window the gain is set by the neighbours,
\begin{equation}\label{eq:learning}
  g_i=r\bigl[(1-\lambda)+\lambda\,\bar s_{V(i)}\bigr]\qquad(a_i\le a_w),
\end{equation}
where $\bar s_{V(i)}$ is the mean competence of the other individuals of the $3\times3$ block around $i$ (zero if
there are none); \Cone\ has $\lambda=1$, purely social learning. A pup starts at $ws_f\le0.2$ and can gain at most
$ra_w=1.2$, but only among competent demonstrators. The rule follows Calhoun's observation that ``by midway in
phase~C essentially all young were prematurely rejected by their mothers''; it is consistent with the non-genomic
transmission of maternal care in rats \citep{Francis1999} and the social learning of pup retrieval in mice
\citep{Carcea2021}, which support the social transmission of maternal care, not of male sexual competence, which
the model extends by hypothesis.
\emph{Crowd aversion (AV) and retreat sites (R7).} The second factor of Eq.~\eqref{eq:move} makes deteriorated
animals ($s\lesssim0.3$) avoid crowded cells, and the third makes them prefer retreat sites, a fraction $f_R$ of the
cells that hold at most $c_R$ individuals. They follow Calhoun's descriptions of withdrawn males that pooled in the
centre, ``beautiful ones'' that never fought, and withdrawn females that retreated to the high, less-preferred nest
boxes \citep{Calhoun1973}.

\begin{table}[!tp]
\centering
\caption{Key parameters of the reference model \Cone\ (SM, Table~\ref{SM-tab:odd-params}, for all parameters and
their justification). The short-lived control differs in $a_{\max}=60$, $a_{\rm rep}=50$, $r=0.01$, $\lambda=0$,
$a_w=\infty$, $w=1$, $\alpha=1.4$ and has neither crowd aversion nor retreat sites.}
\label{tab:params}
\footnotesize
\setlength{\tabcolsep}{3pt}
\begin{adjustbox}{max width=\textwidth}
\begin{tabular}{@{}llp{4.0cm}llp{4.0cm}@{}}
\toprule
Symbol & Value & Meaning & Symbol & Value & Meaning \\
\midrule
$L$ & 30 & side of the arena ($|G|=900$ cells) & $r$ & 0.1 & maximal weekly gain of competence \\
$N_0$ & 200+200 & founders, age 4 weeks, $s=1$ & $a_w$ & 12 weeks & sensitive window (R1) \\
$m_c$ & 8 & tolerated occupancy & $\lambda$ & 1 & weight of social learning (R2) \\
$\gamma$ & 0.1 & damage per extra occupant & $\eadult$ & 0 & adult recovery factor \\
$p_0$ & 0.3 & weekly conception probability & $w$ & 0.2 & maternal transmission \\
$q_0$ & 0.95 & pup survival at $s_f=1$ & $\alpha$ & 4 & attraction to acquaintances \\
litter & 3--10, mean 6.55 & litter size & $\kappa$, $p$ & 0.15, 4 & crowd aversion (AV) \\
$a_{\min}$, $a_{\rm rep}$ & 6, 80 weeks & reproductive ages of females & $f_R$, $c_R$ & 0.3, 2 & retreat sites (R7) \\
$a_{\max}$, $k$ & 120 weeks, 0.15 & senescence & $\pi$, $q$ & 100, 4 & preference for retreat sites (R7) \\
\bottomrule
\end{tabular}
\end{adjustbox}
\end{table}

\paragraph{Order of processes and demography}
Each week, in this order: every individual moves (Eq.~\ref{eq:move}, simultaneously, with the capacity of the
retreat sites enforced); affinities are updated; occupancies are counted; competence is updated
(Eqs.~\ref{eq:stress} and~\ref{eq:learning}); females mate in their cell, conceive and give birth to the pups that
survive, which are placed in the mother's cell; individuals age and die; observables are recorded. The founders,
$200$ females and $200$ males of age 4 weeks with $s=1$, are placed at random. Four simplifications of mouse
demography have measured effects (SM, Fig.~\ref{SM-fig:robustness}). There is no gestation or interbirth interval:
conception and birth occur in the same weekly step, with a ceiling of about two pups per female and week; with a
refractory period of three weeks after each birth \Cone\ still collapses (60 of 60 runs extinct with either of two seed
sets; joint outcome 59 of 60 with an independent seed and 60 of 60 with the production seeds), and the extinction
time changes by $+1.0$ [$-2.6$, $4.6$] weeks (paired with the production seeds). Pups move
and add to crowding from birth; keeping them in the maternal cell for three weeks lowers the joint outcome (below) to
0.83. The founders form a synchronous cohort, and 69\% of all births are theirs; initial ages spread over 4--80 weeks
give 0.63. There is no mortality before senescence; a background mortality of 0.003 per week lowers the joint
outcome to 0.55 and 0.008 to 0.06, in both cases through the timing of the end of natality, while up to 0.05 per
week does not rescue the colony (SM, Section~\ref{SM-supp:recovery}). Longevity is part of the mechanism, not a detail
(Section~\ref{sec:res-mechanism}).

\paragraph{Model structures compared}
The reference model \Cone\ has the core processes and the four social rules. The \emph{core model} keeps the core
processes but replaces the four social rules by neutral values: no window, individual recovery at a slow rate ($r=0.01$) at every age,
full maternal transmission ($w=1$), no crowd aversion and no retreat sites. A null model has no crowding damage
($\gamma=0$). The \emph{short-lived control} (Table~\ref{tab:params}) regenerates competence individually; in space 89
of its 200 colonies are still alive at week $1\,000$, cycling with a period of 162 [155, 170] weeks. The
\emph{reversible-suppression null} is \Cone\ with the two recovery parameters changed: adults recover at a fraction
$\eadult=0.8$ of the juvenile rate, and $\lambda=0.9$ instead of 1 keeps social learning but lets them recover partly
on their own, at the asocial rate $\eadult(1-\lambda)=0.08$ of the juvenile rate. With $\lambda=1$, $\eadult=0.8$
lies inside the transition (10 of 20 colonies die out at $L=30$); the null lies above the rescue threshold at every
size simulated (Section~\ref{sec:res-recovery}). The parameters that set the form of the rules were chosen to reproduce the most
qualitative patterns with the fewest rules; they were not fitted to Calhoun's curves. In the \emph{well-mixed}
variants every individual moves each week to a uniformly random cell, so that attraction, aversion and the
preference for retreat sites have no effect and the retreat sites act only as a capacity; cell occupancies are
Poisson with mean $N/|G|$. A female then mates either within her cell, as in space (\emph{local mating}), or with a
reproductive male drawn from the whole arena (\emph{global mating}), which removes any difficulty of finding a mate.

\subsection{Patterns, evaluation and statistical design}
\label{sec:patterns}

From Calhoun's report we extracted seventeen qualitative, falsifiable patterns and four anti-patterns, each
translated into an operational criterion on simulated time series (SM, Table~\ref{SM-tab:observations-full}; codes
in parentheses). Growth: latency, exponential growth and deceleration (O1--O3, scored only with a separate founder
protocol from four pairs), and a socially functional growth phase (O10b). Peak: a peak far below capacity with empty
nest sites (O4), adults that deteriorate before it (O10), aggregation with free space (O14), and coexisting crowded
and withdrawn non-breeders (O13), with a persistent withdrawn class (O13p) born late (O13b). Decline: late-born
cohorts that fail to mature (O11) and rarely conceive (O11b), natality that ceases near the peak (O5), no recovery of
natality (O6), extinction (O7), a long senescent decline (O8) to an old, female-biased remnant (O9), total infant
mortality (O15), and the failure of transplanted late animals to breed (O12). We call O5, O6, O7 and O10 the
\emph{collapse envelope}. Calhoun's dates enter only through dimensionless ratios of the time and size of the peak.
A criterion returns PASS, MARGINAL (one step from its threshold), FAIL or not applicable for each run, and an
ensemble is summarised by the fraction of passing runs with a 95\% Wilson score interval. The joint
outcome O17p requires every applicable essential pattern, no anti-pattern and a persistent withdrawn class. Every run
was scored with the evaluator preregistered before the production runs (a time-stamped plan in the code repository,
1 October 2026; SM, Section~\ref{SM-supp:evaluators}) and with a primary evaluator, fixed before those runs, that dates the peak at the start of the plateau rather than at $\arg\max N$; the two are reported side
by side (SM, Table~\ref{SM-tab:res-rescoring}). One option, the end of natality in O5 at the week by which 99\% of the
births have occurred, was chosen after exploratory runs had shown that the last sustained birth fails O5 in about
half of the \Cone\ runs (with that reference O17p is 0.485, 97 of 200 production runs). Square brackets give 95\% intervals unless
stated otherwise.

The unit of analysis is the independent stochastic replicate. The design comprises production ensembles of the model
structures ($n=200$ each), a factorial ablation of the rules, transplant experiments ($60$ replicates per arm), a
sweep of the adult recovery factor at three arena sizes, robustness variants, lattice sizes, and a global
sensitivity analysis (Morris screening, two Latin hypercubes and Sobol' indices); seven
confirmatory hypotheses were tested with Holm correction, and extinction times are summarised by Kaplan--Meier
estimates (SM, Sections~\ref{SM-supp:Q} and~\ref{SM-supp:ensembles}).

\paragraph{Use of AI in the research}
Starting from the author's preliminary implementation and draft rules, Claude (Anthropic) agents, used through Claude
Code (v2.1; Claude Opus 5.5 and Claude Fable 5.1 for the research tasks, Claude Sonnet 5.5 for literature searches and
auxiliary programming), reimplemented and extended the code, ran the simulations and the statistical analyses, and
searched the literature. The author set the objectives, the qualitative targets and the interaction rules, decided
every change of scope and approved each stage. The code has an automated test suite; scripts in the code repository
regenerate every figure and number from the archived outputs; independent Claude agents reviewed the work, checking
the code against the text and reproducing key results; and the history is kept under version control.
